% !TeX root = ../main.tex
\section{Type inference and checking}

STLC's typing rules (\S3) are syntax-directed, so given a fully-annotated term, checking it against a type is entirely mechanical. But nobody wants to write $\lambda x{:}\tyint.\,\lambda y{:}\tyint.\,\lambda f{:}\tyint\to\tyint\to\tyint.\,f\,x\,y$ when \lstinline|\x y f -> f x y| conveys the same program. This section is about the gap between those two: how much of the typing can the machine reconstruct on its own, and what exactly is the machine doing when it does?

\subsection{checking vs.\ inference: two different problems}

It's worth pulling these apart carefully, because they really are different problems with different algorithms, even though people say ``type checking'' for both colloquially.

\begin{description}
  \item[Type checking] Given $\Gamma$, $e$, \emph{and a candidate type} $\tau$, decide whether $\Gamma \vdash e : \tau$ holds. You already have a hypothesis; the job is just to verify it. For STLC with fully annotated lambdas, this is exactly what \S3's syntax-directed rules do — read the term, follow the one applicable rule at each step, and either every subgoal succeeds or one fails.
  \item[Type inference] Given $\Gamma$ and $e$ \emph{alone} (no annotations, no candidate type), find a $\tau$ such that $\Gamma \vdash e : \tau$ — or report that none exists. Now there's no hypothesis to check; the type itself has to be discovered from how $e$'s pieces are used.
\end{description}

Checking is the easier problem in general: it's a verification task, and for STLC it terminates trivially by structural recursion on $e$. Inference is strictly harder — you're solving for unknowns rather than confirming a guess — and for sufficiently expressive type systems (System F in full generality, for instance) it's actually undecidable. The classical algorithm for inference in the Hindley-Milner fragment (ML, Haskell's core, OCaml) is \textbf{Algorithm W}, and it works by turning ``guess the type'' into ``solve a system of equations'', which is a problem we know how to attack systematically.

\subsection{Algorithm W: constraints, then unification}

The idea: don't try to guess types directly. Instead, walk the term once, inventing a fresh type variable for anything unknown and recording an equation every time the typing rules force two types to coincide; then solve the whole system of equations at the end. Concretely:

\begin{enumerate}
  \item \textbf{Generate.} Walk $e$ top-down. For each subexpression whose type isn't yet known, assign a fresh type variable $\alpha_i$. For each typing rule that applies, emit a \emph{constraint} $\tau = \tau'$ recording what that rule requires — e.g., \textsc{T-App}'s requirement that the function's domain match the argument's type becomes a constraint $\tau_{\mathrm{fun}} = \tau_{\mathrm{arg}} \to \beta$ for a fresh $\beta$ standing for the result.
  \item \textbf{Solve.} Feed the whole constraint set to \textbf{unification}: find a substitution $\theta$ (an assignment of types to type variables) that makes every constraint hold syntactically. Unification proceeds constraint-by-constraint, decomposing structurally (unifying $\tau_1\to\tau_2$ with $\sigma_1\to\sigma_2$ recurses into unifying $\tau_1$ with $\sigma_1$ and $\tau_2$ with $\sigma_2$) and binding a variable to a type whenever one side is a bare variable — \emph{except} that binding $\alpha$ to $\tau$ is refused if $\alpha$ occurs free in $\tau$ (the \textbf{occurs check}; see the worked example below and \S7's discussion of exactly why this restriction is there).
  \item \textbf{Report.} If unification fails at any point (a structural mismatch, like trying to unify $\tyint$ with $\tybool\to\tybool$, or a failed occurs check), the program has a type error. If it succeeds, apply the resulting substitution $\theta$ to $e$'s type variable, and that's $e$'s inferred type.
\end{enumerate}

\subsection{principal types}

A natural question: when inference succeeds, is the type it produces the ``right'' one, or just \emph{a} type that happens to work? For the Hindley-Milner fragment the answer is clean: unification, when it succeeds, produces a \textbf{most general unifier} — the least commitment substitution that solves the constraints, with every other valid solution obtainable by substituting further into it. This translates directly into a \textbf{principal types} property for the type system: whenever $e$ is well-typed at all, Algorithm W infers a single type $\tau$ that is the \emph{most general} type $e$ can have — every other valid typing of $e$ is an instance of $\tau$ (obtained by substituting concrete types for $\tau$'s free type variables). E.g., $\lambda x.\,x$'s principal type is $\alpha\to\alpha$; $\tyint\to\tyint$ is \emph{a} valid type for it too, but it's a specific instance, not the principal one. This is exactly the type that gets $\forall$-quantified and generalized at a \lstinline|let| in \S5's discussion of let-polymorphism — generalization only makes sense because there's a canonical, most-general type to quantify over.

\subsection{solved example: $\lambda f.\, f\, f$ does not typecheck}

Consider \lstinline|proc(f) zero?((f f))|. We try to infer the type of $f$.

\begin{itemize}
  \item Let $f : \alpha$ (fresh).
  \item The expression $f\, f$ applies $f$ to itself, so \textsc{T-App}'s constraint-generation rule demands $f$'s type be a function type whose domain is (again) $f$'s type: $\alpha = \alpha \to \beta$ for some fresh $\beta$.
  \item Unification tries to bind $\alpha \mapsto \alpha\to\beta$ — but the occurs check catches that $\alpha$ appears free on the right-hand side, and refuses the binding. (Ignoring the check and expanding anyway shows why: $\alpha = (\alpha\to\beta)\to\beta = ((\alpha\to\beta)\to\beta)\to\beta = \cdots$ never bottoms out at a finite type.)
\end{itemize}

Unification fails, so the program has no type in the Hindley-Milner fragment (equivalently, in STLC — this is the same self-application that \S3 showed defeats strong normalization, and the same equation \S7 revisits when it introduces $\mu$-types as the deliberate, explicit way to give equations like $\alpha=\alpha\to\beta$ a solution).

\subsection{typing rules for references}

Inference and checking aren't limited to the pure lambda calculus — the same machinery extends to imperative constructs once we add a store. Reference cells get a type $\tref{T}$ (``a mutable cell holding a $T$''), and dereference, assignment, and \lstinline|while| get straightforward rules:
\[
\infer[\textsc{T-Deref}]
  {\Gamma \vdash \deref{l} : T}
  {\Gamma \vdash l : \tref{T}}
\qquad
\infer[\textsc{T-Assign}]
  {\Gamma \vdash \assign{l}{e} : \unit}
  {\Gamma \vdash l : \tref{T} \quad \Gamma \vdash e : T}
\qquad
\infer[\textsc{T-While}]
  {\Gamma \vdash \whiledo{B}{C} : \unit}
  {\Gamma \vdash B : \tybool \quad \Gamma \vdash C : \unit}
\]
Note that \textsc{T-Assign} and \textsc{T-While} both give type $\unit$ to the whole construct — commands are executed for their effect on the store, not for a value, so $\unit$ is the honest ``there's nothing here'' type. These rules slot into constraint generation exactly like \textsc{T-App} did above: e.g., \textsc{T-Deref} applied to a variable of unknown type generates the constraint that its type be $\tref{\beta}$ for a fresh $\beta$.

\subsection{worked example: infer types for the following program}

Here's a small imperative program using the constructs above; let's read off the type of every piece the way Algorithm W's constraint-generation step would, then check the results are consistent.

\begin{lstlisting}
B  = !n < 5
B1 = !term <= 2
C1 = res := !n
C2 = c:=!a+!b; res:=!n; !c; a:=!b; b:=!c
C  = if B1 then C1 else C2; n:=!n+1
a:=1; b:=1; while B do C
\end{lstlisting}

Reading off types from the operations performed:
\begin{itemize}
  \item $n, \mathit{res}, a, b, c : \tref{\tyint}$ — each is dereferenced with $!$ and used in arithmetic or comparison against an integer literal, so \textsc{T-Deref} forces their referent type to be $\tyint$.
  \item $\mathit{term} : \tref{\tyint}$ — compared with the literal $2$, same reasoning.
  \item $B, B_1 : \tybool$ — both are comparisons (\lstinline|<|, \lstinline|<=|), and \textsc{T-While} requires $B:\tybool$ to typecheck at all.
  \item $C_1, C_2, C : \unit$ — each is a sequence of assignments (and one nested \lstinline|if|), and \textsc{T-Assign}/sequencing both produce $\unit$; \textsc{T-If} then requires both branches ($C_1$ and $C_2$) to agree on type, which they do.
  \item The whole program : $\unit$ — it's a sequence of assignments ending in a \lstinline|while|, and \textsc{T-While} itself has type $\unit$.
\end{itemize}
Every constraint generated along the way (e.g., $n$'s referent type $=\tyint$, from both the comparison in $B$ and the arithmetic in $C_2$) is consistent with every other, so unification succeeds and the whole program typechecks.

\subsection{aside: Scala generics}
\begin{lstlisting}[language=Scala]
case class MyPair[A,B](x:A,y:B)
val p = MyPair(1, "Scala")
def id[T](x:T)=x
val q = id(1)
\end{lstlisting}
The definition of \lstinline|id| introduces a \texttt{polymorphic} type: it can be applied to any type, and each call site independently instantiates a fresh $T$ — this is Scala's version of exactly the generalize-at-binding-site behavior \S5 describes for \lstinline|let|-polymorphism, just spelled with explicit type parameters instead of let-generalization. \lstinline|MyPair| similarly generalizes over two independent type parameters $A,B$, inferred here from the arguments $1$ and \lstinline|"Scala"| without any annotation — Scala's local type inference is doing a restricted version of the same constraint-solving Algorithm W does, though (being a nominally-typed, subtyping-aware language) it falls short of full HM in some corners, which is part of why Scala still asks for annotations more often than OCaml does.

\subsection{bidirectional type checking}

Full Hindley-Milner inference is elegant, but it has a real cost: because the algorithm commits to unifying types eagerly and globally, it interacts poorly with features like subtyping (\S6) and higher-rank polymorphism (\S5) — a single global unification pass doesn't have an obvious place to consult a subtyping relation or to intelligently pick when a $\forall$ should be introduced or eliminated. Most modern typed languages (TypeScript, Rust, Swift, and research languages built on top of bidirectional systems) instead use \textbf{bidirectional type checking}, which sidesteps full inference by splitting the single judgment $\Gamma\vdash e:\tau$ into two modes that call each other recursively:

\begin{description}
  \item[Synthesis (or inference) mode,] written $\Gamma \vdash e \Rightarrow \tau$: given $e$, \emph{compute} its type $\tau$ from the bottom up. This works when $e$'s structure carries enough information on its own — variables (look up $\Gamma$), literals (obvious), function application (synthesize the function's type, then check the argument against its domain).
  \item[Checking mode,] written $\Gamma \vdash e \Leftarrow \tau$: given \emph{both} $e$ and an expected type $\tau$, verify $e$ fits. This handles constructs that don't carry enough local information to synthesize a type on their own — most notably, a bare, unannotated $\lambda x.\,e$: you cannot synthesize $\lambda x.\,e$'s type without knowing what type $x$ should have, but if you're \emph{told} the expected type is $\tau_1\to\tau_2$, checking mode can assume $x:\tau_1$ and check $e \Leftarrow \tau_2$.
\end{description}

The two modes are mutually recursive and mostly syntax-directed (in contrast to full HM, no global unification pass is needed): synthesis handles the ``information-rich'' constructs and hands expected types down into checking mode for the ``information-poor'' ones, and a single rule bridges the two directions — synthesize $e$'s type, then check it matches (up to subtyping, if that's in the system) what was expected:
\[
\infer[\textsc{Sub/Check}]
  {\Gamma \vdash e \Leftarrow \tau}
  {\Gamma \vdash e \Rightarrow \tau' \quad \tau' <: \tau}
\]
This is a big part of why bidirectional checking composes so well with subtyping (\S6) and rank-2+ polymorphism (\S5): both features are about ``does this fit where it's expected'', which is exactly what checking mode is built to ask, without needing a global unification algorithm to know where subtyping or a higher-rank instantiation should kick in. The price, compared to full HM, is that you generally need \emph{more} annotations — a bare lambda passed where no expected type is available (e.g., as the direct argument to an un-annotated generic function) can genuinely fail to synthesize a type, where Algorithm W would have solved for it via unification. Most practical bidirectional systems patch this with local, limited unification for exactly those cases, rather than reverting to full global HM.

\subsection{Exercises}

\begin{itemize}
  \item Infer the type of the S combinator: $\lambda x.\,\lambda y.\,\lambda z.\; x\,z\,(y\,z)$. What constraints does unification generate, and what is the principal type?

  \item Consider $\lambda f.\,\lambda x.\; f\,(f\,x)$. Does it typecheck? If so, give its most general (principal) type.

  \item What types are inferred for \lstinline|f|, \lstinline|result|, and \lstinline|pairs| in the following Scala snippet?
\begin{lstlisting}[language=Scala]
def f[A](lst: List[A]) = lst.head
val result = f(List(1,2,3))
val pairs  = List((1,"a"), (2,"b"))
\end{lstlisting}

  \item Given $x : \tref{\tyint}$ and $b : \tref{\tybool}$, does the following expression typecheck? If so, give its type.
\[
  \assign{b}{\deref{x}=0}\;;\;\ift{\deref{b}}{\assign{x}{1}}{\mathbf{skip}}
\]

  \item Explain, precisely, the difference between ``the program has no type'' (as with $\lambda f.\,f\,f$) and ``the algorithm cannot decide the type without more information'' (as with an unannotated lambda passed to a bidirectional checker with no expected type available). Are these the same kind of failure?

  \item Trace the synthesis/checking modes by hand for $(\lambda x{:}\tyint.\,x)\,3$ under bidirectional type checking: which subterms are synthesized, which are checked, and against what expected type?
\end{itemize}
